\documentclass[../main.tex]{subfiles}
\graphicspath{{\subfix{../image/}}}

\begin{document}

\chapter{线性回归专题}\label{chapter:appE-线性回归专题}

本附录收录关于线性回归模型的两个结果：二次型期望公式与 OLS 误差方差的无偏估计。它们在本书正书中未涉及，但在计量经济学与统计推断中常用，故附于此。

\begin{lemma}[二次型期望公式]\label{lemma:appE-二次型期望公式}
设 $\varepsilon$ 为 $n$ 维随机向量，$E(\varepsilon)=\mu$，$\operatorname{Cov}(\varepsilon)=\Sigma$，$A$ 为 $n\times n$ 实矩阵，则
\begin{align*}
E(\varepsilon^T A \varepsilon) = \operatorname{tr}(A\Sigma) + \mu' A \mu.
\end{align*}
特别地，若 $\mu=0$，$\Sigma=\sigma^2 I_n$，则 $E(\varepsilon^TA\varepsilon)=\sigma^2\operatorname{tr}(A)$。
\end{lemma}

\begin{proof}
设 $\varepsilon =(\varepsilon_1,\cdots,\varepsilon_n)^T$，$\mu =(\mu_1,\cdots,\mu_n)^T$，$\Sigma =(\Sigma_{ij})_{n\times n}$，$A=(a_{ij})_{n\times n}$，则
\begin{align*}
\varepsilon^T A \varepsilon = \sum_{i=1}^n\sum_{j=1}^n a_{ij}\varepsilon_i\varepsilon_j.
\end{align*}
从而
\begin{align}\label{align:appE-二次型期望}
E(\varepsilon^T A\varepsilon )=\sum_{i,j=1}^n{a_{ij}E(\varepsilon _i\varepsilon _j)}.
\end{align}
由协方差定义知
\begin{align*}
\mathrm{Cov(}\varepsilon _i,\varepsilon _j)&=E\left[ \left( \varepsilon _i-E\left( \varepsilon _i \right) \right) \left( \varepsilon _j-E\left( \varepsilon _j \right) \right) \right]
=E\left[ \left( \varepsilon _i-\mu _i \right) \left( \varepsilon _j-\mu _j \right) \right] \\
&=E\left( \varepsilon _i\varepsilon _j \right) -\mu _jE\left( \varepsilon _i \right) -\mu _iE\left( \varepsilon _j \right) +\mu _i\mu _j=E\left( \varepsilon _i\varepsilon _j \right) -\mu _j\mu _i-\mu _i\mu _j+\mu _i\mu _j \\
&=E\left( \varepsilon _i\varepsilon _j \right) -\mu _i\mu _j.
\end{align*}
进而
\begin{align*}
E\left( \varepsilon _i\varepsilon _j \right) =\mathrm{Cov}\left( \varepsilon _i,\varepsilon _j \right) +\mu _i\mu _j.
\end{align*}
将上式代入\eqref{align:appE-二次型期望}式得
\begin{align*}
E(\varepsilon^T A\varepsilon )&=\sum_{i,j=1}^n{a_{ij}\Sigma _{ij}}+\sum_{i,j=1}^n{a_{ij}\mu _i\mu _j}=\sum_{i,j=1}^n{a_{ij}\Sigma _{ji}}+\mu ' A\mu \\
&=\sum_{i=1}^n{(A\Sigma )_{ii}}+\mu' A\mu =\mathrm{tr(}A\Sigma )+\mu' A\mu .
\end{align*}
特别地，若 $\mu=0$，$\Sigma=\sigma^2 I_n$，代入上式得 $E(\varepsilon^TA\varepsilon)=\sigma^2\operatorname{tr}(A)$。
\end{proof}

\begin{theorem}[OLS 误差方差的无偏估计]\label{theorem:appE-OLS误差方差无偏估计}
考虑线性回归模型
\begin{align*}
Y = X\beta + \varepsilon,
\end{align*}
其中：
\begin{itemize}
\item $Y$ 是 $n$ 维观测向量，
\item $X$ 是 $n \times (K+1)$ 列满秩设计矩阵$(K+1<n)$，
\item $\beta$ 是 $(K+1)$ 维未知参数向量，
\item $\varepsilon$ 是 $n$ 维随机误差向量。
\end{itemize}

假设：
\begin{enumerate}[(1)]
\item $E(\varepsilon)=0$，
\item $\operatorname{Cov}(\varepsilon)=E(\varepsilon\varepsilon^T)=\sigma^2 I_n$，其中 $\sigma^2=\operatorname{Var}(\varepsilon_i)>0\ (i=1,2,\cdots,n)$ 是未知常数。
\end{enumerate}

定义：
\begin{itemize}
\item OLS 估计量：$\hat{\beta}=(X^TX)^{-1}X^TY$，
\item 残差向量：$e=Y-X\hat{\beta}$，
\item 残差平方和：$e^Te$，
\item 误差方差估计量：$s_e^2=\dfrac{e^Te}{n-K-1}$。
\end{itemize}

则 $s_e^2$ 是 $\sigma^2$ 的无偏估计，即
\begin{align*}
E(s_e^2)=\sigma^2.
\end{align*}
\end{theorem}

\begin{proof}
由题设可知
\begin{align*}
e &= Y - X\hat{\beta} = (X\beta+\varepsilon) - X[(X^TX)^{-1}X^T(X\beta+\varepsilon)] \\
&= X\beta+\varepsilon - X\beta - X(X^TX)^{-1}X^T\varepsilon \\
&= \bigl[I_n - X(X^TX)^{-1}X^T\bigr]\varepsilon.
\end{align*}
记 $M = I_n - X(X^TX)^{-1}X^T$，则 $e=M\varepsilon$。显然 $M$ 是对称幂等阵，从而
\begin{align*}
e^Te = (M\varepsilon)^T(M\varepsilon)=\varepsilon^T M' M \varepsilon = \varepsilon^T M \varepsilon,
\end{align*}
由假设知 $\varepsilon$ 满足 $E(\varepsilon)=0$，$\operatorname{Cov}(\varepsilon)=\sigma^2 I_n$，则由\cref{lemma:appE-二次型期望公式}得
\begin{align*}
E\left( e^Te \right) &=\sigma ^2\mathrm{tr}\left( M \right) =\sigma ^2\mathrm{tr}\bigl( I_n-X\left( X^TX \right) ^{-1}X^T \bigr) \\
&=\sigma ^2\mathrm{tr}\left( I_n \right) -\sigma ^2\mathrm{tr}\bigl( X\left( X^TX \right) ^{-1}X^T \bigr) =\sigma ^2n-\sigma ^2\mathrm{tr}\bigl( \left( X^TX \right) ^{-1}X^TX \bigr) \\
&=\sigma ^2n-\sigma ^2\mathrm{tr}\left( I_{K+1} \right) =\sigma ^2\left( n-K-1 \right) .
\end{align*}
因此
\begin{align*}
E(e^Te)=\sigma^2\bigl(n-(K+1)\bigr)=\sigma^2(n-K-1).
\end{align*}
故
\begin{align*}
E(s_e^2)=E\left(\frac{e^Te}{n-K-1}\right)=\frac{1}{n-K-1}E(e^Te)=\frac{1}{n-K-1}\sigma^2(n-K-1)=\sigma^2.
\end{align*}
\end{proof}

\end{document}